\emph{Quantitative orbit rigidity of the moment coordinate near coincidence walls.} Within the named coupled gaussian slot law at strictly positive profile and the moment coordinates \(\rho_{p}=p!\,h_{p}(g)/W^{p}=2^{p-1}m_{2p}/m_{2}^{p}\) with \(g_{1}+\cdots+g_{D}=W\) of the preceding paragraph, the map from the trace-fixed chamber \(g_{1}>\cdots>g_{D}>0\), with \(g_{D}=W-\sum_{i<D}g_{i}\) eliminated, to \((\rho_{2},\ldots,\rho_{D})\) has in the chamber coordinates \((g_{1},\ldots,g_{D-1})\) the exact Jacobian determinant \(\pm(\prod_{k=1}^{D}k!)\,V(g)/W^{D(D+1)/2-1}\), where \(V(g)=\prod_{i<j}(g_{i}-g_{j})\) is the Vandermonde of the preceding paragraph and the sign is fixed by the ordering convention, verified symbolically at \(D=2\) and \(D=3\) and at three rational profiles at \(D=4\): it vanishes to first order exactly on the coincidence walls and nowhere else in the chamber, so the moment tuple is a genuine chamber coordinate off the walls. At \(D=2\) the fold is exact and global on the chamber: \(\rho_{2}=1+s_{2}/W^{2}=(3W^{2}+(g_{1}-g_{2})^{2})/(2W^{2})\), so the wall value is the uniform-profile value \(3/2\), equal to the uniform lower bound \(\prod_{k=1}^{p-1}(1+k/D)\) at \(p=2\) and \(D=2\), the single-row limit value is \(2=p!\), and the separation obeys the exact inversion \((g_{1}-g_{2})^{2}=2W^{2}(\rho_{2}-3/2)\); the inverse moment map is globally \(1/2\)-H\"older with constant \(\sqrt{2}\,W\), this exponent is exact at the wall, and on \(\lvert g_{1}-g_{2}\rvert\ge\delta\) the inverse is Lipschitz with constant at most \(W^{2}/\delta\). At general \(D\), along the wall \(g_{1}=g_{2}=c\) with the pair mean and the remaining slots fixed, the power-sum transverse deviation is the exactly even polynomial \(s_{k}(c+u)+s_{k}(c-u)-2s_{k}(c)=2\sum_{j\ge1}\binom{k}{2j}c^{k-2j}u^{2j}\) with no odd terms, the \(\rho_{2}\) deviation is exactly \((2/W^{2})u^{2}\) independently of the wall position, and the scaled power-sum deviations \(\sigma_{k}=s_{k}(g)/W^{k}\) have leading transverse coefficients \(k(k-1)c^{k-2}u^{2}/W^{k}\), so the fold \(u\mapsto u^{2}\) is the transverse normal form at simple walls; the smallest singular value of the coordinate differential vanishes like the gap at a single wall and like \(\delta^{D-1}\) along profiles whose gaps all have order \(\delta\), and the sharp inverse modulus is \(1/2\)-H\"older also along the symmetric path into the triple point, these rates being checked at \(D=3\). On the named fixed-trace Dirichlet-beta family, on the \(D=2\) Gram separation coordinate \(u=\lvert p_{1}-p_{2}\rvert\) with density proportional to \(u^{\beta}(1-u^{2})^{a-1}\), the moment \(\mathrm{E}\,\lvert p_{1}-p_{2}\rvert^{-1}\) at the computed \((a,\beta)=(\lvert N-2M\rvert+1,2)\) member is exactly \((2a+1)\Gamma(a+\tfrac12)/(a\sqrt{\pi}\,\Gamma(a))\), equal to \(3/2\) at \(a=1\), while the \(\beta=0\) members of the family diverge logarithmically with leading coefficient \(2\Gamma(a+\tfrac12)/(\sqrt{\pi}\,\Gamma(a))\), equal to \(1\) at \(a=1\); against these weights \(\int\lvert V\rvert^{-1}\,d\mu\) converges near an \(m\)-fold coincidence exactly for \(\beta>1-2/m\), in particular \(\beta>0\) at a simple wall and \(\beta>1/3\) at the triple point of \(D=3\), so within the named family the inverse-Jacobian moment is finite exactly when the interaction exponent exceeds the coincidence threshold. These statements are exact change-of-variables, normal-form and beta-integral statements within the named coupled gaussian slot law at strictly positive profile with the moment convention of the preceding paragraph, and within the named fixed-trace Dirichlet-beta family on the Gram coordinate; the moment tuple identifies orbits, with the two chamber points over a wall related by square-root branching, so it selects no chamber point, no profile and no group element, they evaluate no matrix integral of the core action, they fix no state, they leave the missing steps (i)-(iv) of the Missing step paragraph unchanged, and they close nothing in Problem~G.
